\section{几何替换必须支付什么}\label{sec:replacement-bank}
以下比较的实际并集均可测且面积有限。
设一个构造保留背景 $B$，撤去一批并集为 $P$ 的三角形，再插入新的完整单位针，其三角形的并集为 $Q$。面积变化恰好是
\begin{equation}\label{eq:physical-change}
 |B\cup Q|-|B\cup P|=|Q\setminus B|-|P\setminus B|.
\end{equation}
把每个并集分成保留的背景及其外部部分，就得到这个恒等式。这个恒等式也能用来诊断替换。一个三角形可能缩小，但它原来占用的材料被许多其他三角形重复覆盖，而替换物落到了所有其他三角形外。此时把各来源的节省相加，算出的既不是式~\eqref{eq:physical-change} 的左边，也不是右边。

对紧致原点星形体，同一套面积核算可以写成极坐标形式。若 $\rho_K(\phi)$ 是它在实际射线 $\phi$ 上的径向延伸长度，则
\[
 |K|=\frac12\int_0^{2\pi}\rho_K(\phi)^2\,d\phi.
\]
所有记录都在定义 $\rho_K$ 的同一个径向上确界中竞争。来源方向 $q$ 与实际射线 $\phi$ 的作用不同：一个测度为零的方向标签仍可能贡献一个非退化三角形，而一条记录可以占据一整个实际射线区间。下面的例子会保留这些三角形，而不随标签一起将它们丢掉。

成功的上界构造恰好利用了这种竞争。它们的纵向分拆同时移动两个端点的相位，而有符号高度把完整底边的部分放到已经存在的三角形后面。随后的高度和纵向变分移去暴露的条带，同时把补偿材料留在已有的径向延伸范围内。最后的弦长延长步骤使用体内实际包含的整个区间，而不只是两个可见端点。一致的删去量在有限恢复之前就已得到保证。显式构造及其严格比较需要各自的几何假设；替换恒等式解释了这些比较为什么有意义，却不提供那些假设。

